\section{Exercise 5}
\textit{Consider the function defined $f(x,y) = xy+14$ and let
$g(x,y) = x^2+y^2-18$.}
\begin{itemize}
    \item Sketch the level-set for $f$: $L_1 = \{(x,y) \in | f(x,y) = 1\}$ and level-set for h: $M_0 = \{(x,y) \in | g(x,y) = 0\}$ in the same coordinate system.
    \item Find the maximum and minimum for $f$ subject to the constraint $g(x,y) = 0$ (argue for your calculations).
\end{itemize}
%A
the sketch of the general shape of the function $L_1$ and $M_0$, can be done by selecting a few points which satisfies the equation
for each of the sets. In $f(x,y) = xy+14$ the level is for $f(x,y)=1$, and since there a $+14$ in it can be solved such that:
\[f(x,y) = xy+14 \Rightarrow xy+14=1 \Rightarrow xy=-13, \qquad x\neq 0, \qquad y\neq 0\]

\begin{table}[H]
    \centering
    \begin{tabularx}{\textwidth}{|X|X|X|}
        \hline
        $x$ & $y$ & $f(x,y)$ \\ \hline
        -13  & 1   & 1 \\ \hline
        -6.5 & 2   & 1 \\ \hline
        -1   & 13  & 1 \\ \hline
        1    & -13 & 1 \\ \hline
        6.5  & -2  & 1 \\ \hline
        13   & -1  & 1 \\ \hline
    \end{tabularx}
    \caption{Points belonging to the level set $L_1$.}
    \label{tab:levelset}
\end{table}
%B
Since $f$ and $g$ are functions with continuous first-order partial derivatives, this can solved using lagrange multipliers.
Since we're trying to find the minimum and maximum of a function, it is no longer possible to choose any combination of x and y,
since the condition for $g(x,y) = 0$ needs to hold.
Starting with the gradients:
\[\nabla f(x,y) = 
\begin{bmatrix}
    \frac{\partial f}{\partial x}\\
    \frac{\partial f}{\partial y}
\end{bmatrix}
=
\begin{bmatrix}
    y\\x
\end{bmatrix}
,\qquad
\nabla g(x,y) =  
\begin{bmatrix}
    \frac{\partial g}{\partial x}\\
    \frac{\partial g}{\partial y}
\end{bmatrix}
=
\begin{bmatrix}
    2x\\2y
\end{bmatrix}\]
Yielding two equations and the $g(x,y)$ function with the constraint
\[
\begin{cases}
    y = \lambda \ 2x\\
    x = \lambda \ 2y\\
    x^2 + y^2 = 18
\end{cases}
\Rightarrow
\begin{cases}
    \lambda = \frac{y}{2x}\\
    \lambda = \frac{x}{2y}\\
    x^2+y^2=18
\end{cases}
\Rightarrow
\begin{cases}
    \frac{y}{2x} = \frac{x}{2y}\\
    x^2+y^2=18
\end{cases}
\Rightarrow
\begin{cases}
    \frac{y}{x} = \frac{x}{y}\\
    x^2+y^2=18
\end{cases}
\]
\[
\begin{cases}
    y^2 = x^2\\
    x^2+y^2=18
\end{cases}
\Rightarrow
\begin{cases}
    y^2 = \pm x\\
    2x^2 = 18
\end{cases}
\Rightarrow
    x = \pm \sqrt{\frac{18}{2}} = \begin{bmatrix}
        3\\-3
    \end{bmatrix}
\]
Since it has been shown that $y^2=x^2 \Rightarrow y = \begin{bmatrix}3\\-3\end{bmatrix}$
Maximum and minimums is then determined by putting those x and y values into $f(x,y)$
\[f(x,y) \Rightarrow 
\begin{cases}
    f(3,3) = 23\\
    f(3,-3) = 5\\
    f(-3,3) = 5\\
    f(-3,-3) = 23
\end{cases}\]
\[\]
